% Full arbitrary-order partition negation and Stavi completeness.

\subsection{The full arbitrary-order Stavi completeness construction}
\label{sec:p09-stavi-completeness}

We complete the partition-negation argument for the explicitly
defined basis \(\mathcal T_{\mathrm{gap}}\). The labels below
are supplied finite formulas of that language, interpreted
in the original whole model. The construction works on
every nonempty strict linear order; no order-completeness,
density or discreteness assumption is introduced
\cite[Propositions~5.3--5.4 and Theorem~4.4]{Rabinovich2018Stavi}.

\begin{hbdefinition}{def:p09-stavi-completion-primitives}{The elementary interval conditions}
Throughout a bounded construction assume actual \(a<b\).
Write \(E(a,b)\) for emptiness of \((a,b)\),
\(\operatorname{Occ}_B(a,b)\) for
\(\exists z(a<z<b\land B(z))\), and
\(\operatorname{All}_B(a,b)\) for
\(\forall z(a<z<b\to B(z))\).
For a tuple with flags \(W=(\langle D_1,\ldots,D_N\rangle,O)\),
write \(\Pi_W(a,b)\) for its partition of \((a,b)\).
Other bracket choices on \(\Pi_W\) indicate
the corresponding closed or half-open interval.
Its subtuple \(\Pi_{i:j}\) uses flags
\[
                     O_{i:j}=\{h:1\leq h\leq j-i+1,\ h+i-1\in O\}.
\]
For \(i>j\), \(\Pi_{i:j}(a,b)\) is the bookkeeping
condition \(E(a,b)\), not a positive-length partition.
Its complement under \(a<b\) is
\(\operatorname{Occ}_{\top}(a,b)\).
\end{hbdefinition}

\begin{hblemma}{thm:p09-stavi-completion-primitives}{proved=annex:p9.stavi.completion.primitives}{Simple conditions and exhaustive point-or-gap boundaries}
The elementary conditions and their negations have
effective simple partition presentations:
\[
\begin{aligned}
 E(a,b)&=\operatorname{Part}(\langle\top,\top\rangle,\{1,2\})[a,b],\\
 \operatorname{All}_B(a,b)&=
       E(a,b)\lor\operatorname{Part}(\langle B\rangle,\varnothing)(a,b),\\
 \neg\operatorname{All}_B&=\operatorname{Occ}_{\neg B},\qquad
 \neg\operatorname{Occ}_B=\operatorname{All}_{\neg B}.
\end{aligned}
\]
Exact finite cardinalities and their complements are
also simple, using forced singleton cells and increasing
tuples of true-labelled points.

A nonempty occurrence set in \((a,b)\) has its lower
boundary either at \(a\), at an actual point \(c\in(a,b)\),
or at a proper gap \(g\) with \(a<g<b\).
The first case is expressed by \(\partial^+B(a)\);
the actual-point descriptor is
\[
 \operatorname{Inf}_B(a,c,b)=
     a<c<b\land\operatorname{All}_{\neg B}(a,c)
                          \land(B(c)\lor\partial^+B(c));
\]
the proper-gap case is
\(\gamma^+(\neg B)(a)\land\operatorname{Occ}_B(a,b)\).
There are exact order-dual upper-boundary formulas,
using \(\partial^-\), \(\gamma^-\) and
\[
 \operatorname{Sup}_B(a,c,b)=
     a<c<b\land\operatorname{All}_{\neg B}(c,b)
                          \land(B(c)\lor\partial^-B(c)).
\]
The point and gap possibilities are exhaustive
even when an infimum or supremum is missing from
the domain.
\end{hblemma}

\begin{hbdefinition}{def:p09-stavi-point-chain-recursion}{The plain labelled-chain complement}
Put
\[
 A_m(\mathbf P;a,b)=
       \exists x_1\cdots x_m
          (a<x_1<\cdots<x_m<b\land\bigwedge_iP_i(x_i)).
\]
Construct \(O_m\) for its negation under \(a<b\).
At \(m=1\), let \(O_1=\operatorname{All}_{\neg P_1}\).
For \(m\geq2\) retain the following alternatives:
\[
\begin{gathered}
 \operatorname{All}_{\neg P_1}(a,b)\lor
                  \operatorname{All}_{\neg P_m}(a,b),\\
 \partial^+P_1(a)\land O_{m-1}(P_2,\ldots,P_m;a,b),\\
 \partial^-P_m(b)\land O_{m-1}(P_1,\ldots,P_{m-1};a,b),\\
 \exists c\,(\operatorname{Inf}_{P_1}(a,c,b)\land
                              O_{m-1}(P_2,\ldots,P_m;c,b)),\\
 \exists c\,(\operatorname{Sup}_{P_m}(a,c,b)\land
                              O_{m-1}(P_1,\ldots,P_{m-1};a,c)).
\end{gathered}
\]
For the remaining two-gap cases put
\[
\begin{split}
 G_{\mathbf P}(a,b)={}&
 \gamma^+(\neg P_1)(a)\land\operatorname{Occ}_{P_1}(a,b)\\
 &{}\land\gamma^-(\neg P_m)(b)\land\operatorname{Occ}_{P_m}(a,b),\\
 B_{\mathbf P}(a,z,b)={}&
     a<z<b\land\operatorname{Occ}_{P_1}(a,z)
                         \land\operatorname{Occ}_{P_m}(z,b).
\end{split}
\]
Add
\[
 G_{\mathbf P}(a,b)\land
        \operatorname{Part}(\langle\neg P_1,\neg P_m\rangle,\varnothing)(a,b)
\]
and \(G_{\mathbf P}(a,b)\land\exists z(B_{\mathbf P}(a,z,b)\land C_{\mathbf P})\),
where
\[
\begin{split}
 C_{\mathbf P}={}&
 \bigwedge_{i=1}^{m-1}
   \bigl(O_i(P_1,\ldots,P_i;a,z)\lor
                O_{m-i}(P_{i+1},\ldots,P_m;z,b)\bigr)\\
 &{}\land\bigwedge_{i=2}^{m-1}
   \bigl(O_{i-1}(P_1,\ldots,P_{i-1};a,z)
          \lor\neg P_i(z)
          \lor O_{m-i}(P_{i+1},\ldots,P_m;z,b)\bigr).
\end{split}
\]
Empty conjunctions are true. Every recursive chain
has fewer labels. Compile the finite disjunction
using the positive closure of
\cref{thm:p09-stavi-positive-partitions}.
\end{hbdefinition}

\begin{hbtheorem}{thm:p09-stavi-point-chain-complement}{proved=annex:p9.stavi.point.chain.complement}{The complete plain-chain negation construction}
The recursion terminates and constructs a simple
formula \(O_m\) equivalent to \(\neg A_m\)
under \(a<b\), on every linear order.
For the two proper gaps, the first added
partition expresses that the lower gap is
at or after the upper gap; the last
alternative handles their strict opposite
order. All labels, endpoint equalities
at actual bounds, and internal witness
coincidences in the cut are retained.
This proves the source's Proposition~5.4
with the full remaining label lists.
\end{hbtheorem}

\begin{hblemma}{thm:p09-stavi-bounded-sampling}{proved=annex:p9.stavi.bounded.sampling}{The full bounded sampling lemma}
Let \(F_1,\ldots,F_k\) be the finite-seed
formulas of \cref{def:p09-stavi-prefix-recursion}.
For actual \(a_1\leq d\), the following
are equivalent:
\[
\begin{gathered}
 \exists a_2\cdots a_k\,
       (a_1<a_2<\cdots<a_k\leq d
                         \land\bigwedge_{i=1}^kF_i(a_i)),\\
 \exists b\,(a_1\leq b\leq d\land
                 \operatorname{Part}(\langle D_1,\ldots,D_k\rangle,O)[a_1,b]).
\end{gathered}
\]
The proof preserves \(D_i\land F_i\)
on every constructed cell and supplies
all four singleton-flag cases, including
each possible location of the truncation
point. It holds without an attainment
or density assumption.
\end{hblemma}

\begin{hbdefinition}{def:p09-stavi-uniform-handlers}{Independent bounded-prefix and uniform-cell handlers}
For \(W=(\mathbf D,O)\), let \(L_W(a,b)\)
be the simple complement of
\(\exists z(a<z<b\land\Pi_W(a,z])\).
Construct it by padding a true initial
singleton: use \((\top,D_1,\ldots,D_N)\)
with flags \(\{1\}\cup\{i+1:i\in O\}\)
in the finite-prefix recurrence, obtaining
\(F_1,\ldots,F_{N+1}\), and return
\[
              L_W(a,b)=\neg F_1(a)\lor O_N(F_2,\ldots,F_{N+1};a,b).
\]
Let \(R_W(a,b)\) be the exact order dual,
the complement of
\(\exists z(a<z<b\land\Pi_W[z,b))\).
The tuple, flags and all label formulas
are reversed together, and the output
is reversed back.
\end{hbdefinition}

\begin{hblemma}{thm:p09-stavi-uniform-handlers}{proved=annex:p9.stavi.uniform.handlers}{The bounded interfaces and their precise flag hypotheses}
The displayed handlers have their stated
meanings. Moreover,
\[
\begin{aligned}
 1\notin O,\quad \operatorname{All}_{D_1}(a,b)
 &\ \Longrightarrow\
             (\neg\Pi_W(a,b)\leftrightarrow R_W(a,b)),\\
 N\notin O,\quad \operatorname{All}_{D_N}(a,b)
 &\ \Longrightarrow\
             (\neg\Pi_W(a,b)\leftrightarrow L_W(a,b)).
\end{aligned}
\]
They are constructed from plain-chain negation,
not from a same-rank partition-complement
hypothesis. The unflagged endpoint-cell
conditions are essential.
\end{hblemma}

\begin{hbdefinition}{def:p09-stavi-partition-cut}{The full point cut of a partition}
For \(N\geq2\) and \(a<z<b\), define
\[
 A_i(a,z,b)=
 \begin{cases}
  \Pi_{1:i-1}(a,z)\land D_i(z)\land\Pi_{i+1:N}(z,b),&i\in O,\\
  (\Pi_{1:i-1}\lor\Pi_{1:i})(a,z)\land D_i(z)
       \land(\Pi_{i+1:N}\lor\Pi_{i:N})(z,b),&i\notin O.
 \end{cases}
\]
The empty subtuples have the \(E\) meaning
declared above. For an interior index
\(2\leq i<N\), let \(\overline A_i\)
be its explicit De Morgan complement:
replace each nonempty subtuple by its
already constructed smaller-rank complement,
replace \(E\) by \(\operatorname{Occ}_{\top}\),
and replace \(D_i\) by \(\neg D_i\).
In the unflagged case the two alternatives
on each side become a conjunction of
their complements. Put
\[
                       M_W(a,z,b)=\bigwedge_{i=2}^{N-1}\overline A_i(a,z,b).
\]
\end{hbdefinition}

\begin{hblemma}{thm:p09-stavi-partition-cut}{proved=annex:p9.stavi.partition.cut}{Cut equivalence at the same actual point}
For every fixed actual \(a<z<b\),
\[
                      \Pi_W(a,b)\longleftrightarrow\bigvee_{i=1}^N A_i(a,z,b).
\]
A flagged cut cell is exactly \(\{z\}\);
an unflagged one may have a nonempty
portion on either side, or neither.
Consequently any retained condition
\(H(a,z,b)\) can be conjoined with
the complement by replacing it at
that same \(z\) with
\(\bigwedge_i\neg A_i\).
Each interior \(\overline A_i\)
uses only ranks below \(N\).
\end{hblemma}

\begin{hbdefinition}{def:p09-stavi-partition-negation}{The full open-partition negation algorithm}
Construct \(N_W(a,b)\) under the guard
\(a<b\), by induction on \(N\).
Every returned expression is conjoined with
that guard, including each recursive call.
For \(N=1\), if \(1\notin O\), return
\(E(a,b)\lor\operatorname{Occ}_{\neg D_1}(a,b)\).
If \(1\in O\), return emptiness, at
least two interior points, or an
interior failure of \(D_1\), explicitly
\[
 E(a,b)\lor\exists u\,\exists v\,(a<u<v<b)
                      \lor\operatorname{Occ}_{\neg D_1}(a,b).
\]

For \(N\geq2\), first remove forced
endpoint cells. If \(1\in O\), return
\[
 \neg X\top(a)\lor E(a,b)\ \lor\
 \exists z\bigl(a<z<b\land E(a,z)\land
                  (\neg D_1(z)\lor N_{2:N}(z,b))\bigr).
\]
If \(1\notin O\) and \(N\in O\),
return its explicit past counterpart
\[
 \neg Y\top(b)\lor E(a,b)\ \lor\
 \exists z\bigl(a<z<b\land E(z,b)\land
                  (\neg D_N(z)\lor N_{1:N-1}(a,z))\bigr).
\]
Here \(X\top=\bot\mathbin{\mathsf U_s}\top\)
and \(Y\top=\bot\mathbin{\mathsf S_s}\top\).
All recursive partitions have fewer cells.

It remains that \(1,N\notin O\).
Include \(E(a,b)\). On a nonempty
open interval include the following
uniform and endpoint-limit alternatives:
\[
\begin{gathered}
 \operatorname{All}_{D_1}(a,b)\land R_W(a,b),\qquad
 \operatorname{All}_{D_N}(a,b)\land L_W(a,b),\\
 \partial^+\neg D_1(a)\lor\partial^-\neg D_N(b).
\end{gathered}
\]
For actual boundary points put
\[
\begin{aligned}
 I_W(a,z,b)&=\operatorname{Inf}_{\neg D_1}(a,z,b),\\
 J_W(a,z,b)&=\operatorname{Sup}_{\neg D_N}(a,z,b).
\end{aligned}
\]
Include
\[
\begin{split}
 \exists z\bigl(I_W\land M_W
  &{}\land(\neg D_1(z)\lor N_{2:N}(z,b))\\
  &{}\land[
    \neg D_N(z)\lor\operatorname{Occ}_{\neg D_N}(z,b)
         \lor(N_{1:N-1}(a,z)\land R_W(a,z))]\bigr)
\end{split}
\]
and
\[
\begin{split}
 \exists z\bigl(J_W\land M_W
  &{}\land(\neg D_N(z)\lor N_{1:N-1}(a,z))\\
  &{}\land[
    \neg D_1(z)\lor\operatorname{Occ}_{\neg D_1}(a,z)
         \lor(N_{2:N}(z,b)\land L_W(z,b))]\bigr).
\end{split}
\]
All occurrences of \(I_W,J_W,M_W\)
in these displays have the same
arguments \(a,z,b\).

For two proper failure gaps, define
\[
\begin{split}
 G_W(a,b)={}&
 \gamma^+(D_1)(a)\land\operatorname{Occ}_{\neg D_1}(a,b)\\
 &{}\land\gamma^-(D_N)(b)\land\operatorname{Occ}_{\neg D_N}(a,b),\\
 V_W(a,z,b)={}&G_W\land a<z<b
       \land\operatorname{All}_{D_1}(a,z)\land D_1(z)
       \land D_N(z)\land\operatorname{All}_{D_N}(z,b),\\
 B_W(a,z,b)={}&G_W\land a<z<b
       \land\operatorname{Occ}_{\neg D_1}(a,z)
       \land\operatorname{Occ}_{\neg D_N}(z,b).
\end{split}
\]
Include the two alternatives
\[
\begin{split}
 \exists z\bigl(V_W\land M_W
       &{}\land N_{2:N}(z,b)\land L_W(z,b)\\
       &{}\land N_{1:N-1}(a,z)\land R_W(a,z)\bigr),\\
 &\exists z\,(B_W\land M_W).
\end{split}
\]
Their meanings are respectively \(d<z<c\)
and \(c<z<d\), where \(c\) is the first
failure gap of \(D_1\) and \(d\) the last
failure gap of \(D_N\).

The final equal-gap alternative is defined
in the next display. Every output uses
finite positive assembly; no unresolved
negation of a partition remains.
\end{hbdefinition}

\begin{hbdefinition}{def:p09-stavi-equal-gap-branch}{The named common gap and the exact final-label test}
Keep the same \(W\) with unflagged
first and last cells. Define
\[
\begin{aligned}
 T_W&=D_1\mathbin{\mathsf U_s}
                   (D_1\land\operatorname{UG}(D_1,D_N)),\\
 Q_W(a,b)&=
 G_W(a,b)\land
 \operatorname{Part}(\langle D_1,D_N\rangle,\varnothing)(a,b)
                                                  \land\neg T_W(a).
\end{aligned}
\]
The labels in the common-gap test are the
actual first and final labels \(D_1,D_N\).
For \(1\leq i\leq N\) let
\[
 \lambda_i=\operatorname{OG}_{D_1}
                        (\langle D_1,\ldots,D_i\rangle,O_{1:i}),
\]
and let \(\rho_j\), \(1\leq j\leq N\),
be the past dual of the open-gap translator
for \((D_j,\ldots,D_N)\), using the
gap defined from \(b\) by \(D_N\).
Explicitly reverse the tuple, dualize
its labels, replace every flag \(\ell\geq j\)
by \(N+1-\ell\), use
\(\operatorname{OG}_{D_N^{\mathrm d}}\),
then dualize the resulting formula.

Set
\[
 Z_W(a,b)=
 \bigwedge_{i=1}^{N-1}
              (\neg\lambda_i(a)\lor\neg\rho_{i+1}(b))
 \land
 \bigwedge_{i\notin O}
              (\neg\lambda_i(a)\lor\neg\rho_i(b)),
\]
where the second conjunction ranges over
\(1\leq i\leq N\). Include
\(Q_W(a,b)\land Z_W(a,b)\) as the
last branch of \(N_W\).
\end{hbdefinition}

\begin{hblemma}{thm:p09-stavi-equal-gap-branch}{proved=annex:p9.stavi.equal.gap.branch}{The common-gap case without a point at the gap}
\(Q_W(a,b)\) holds exactly when the two
failure gaps in \(G_W\) coincide. Under
this condition the \(\lambda_i,\rho_j\)
refer to the same gap, and
\[
                              \neg\Pi_W(a,b)\longleftrightarrow Z_W(a,b).
\]
The first conjunction complements all
between-cell cuts; the second complements
all cuts through an unflagged cell.
A singleton cell cannot cross a proper
gap. No predicate is evaluated at that
missing point.
\end{hblemma}

\begin{hbtheorem}{thm:p09-stavi-open-partition-complement}{proved=annex:p9.stavi.open.partition.complement}{Full canonical negation of open partitions}
The algorithm constructs a finite simple
formula \(N_W\) satisfying
\[
                         a<b\land\neg\Pi_W(a,b)\longleftrightarrow N_W(a,b)
\]
on every nonempty linear order.
The positive-length partition is false
when \(a\geq b\), so an unguarded
complement adds those order cases.

Every recursive partition call has
strictly fewer cells. Same-rank uniform
subintervals use the independently proved
\(L_W,R_W\) handlers. The empty,
forced-singleton, actual-boundary,
ordered-gap and equal-gap cases are
exhaustive. This completes
Proposition~5.3 in the declared basis.
\end{hbtheorem}

\begin{hbtheorem}{thm:p09-stavi-simple-negation}{proved=annex:p9.stavi.simple.negation}{All simple partition formulas are closed under negation}
The open-partition result yields effective
negation of every basic and every simple
partition formula in the canonical
\(\mathcal T_{\mathrm{gap}}\) expansion.
For a closed interval, enumerate coincident
or strictly ordered endpoints, retain
their labels, and enumerate whether the
first and last cells have nonempty
remainders after removing the endpoints.
Use \(N_W\) on every resulting open
subtuple, with \(E\) for an empty tuple.

Unary ray formulas and whole-order
partition sentences use their completed
temporal translations and Boolean
negation in the canonical label language.
Finite De Morgan distribution and positive
assembly finish the construction.
\end{hbtheorem}

\begin{hbtheorem}{thm:p09-stavi-completeness}{proved=annex:p9.stavi.completeness}{The full past/future Stavi theorem on arbitrary linear orders}
Every monadic first-order formula with
one free variable has an effectively
constructed equivalent formula of
\[
              TL(\mathsf U_s,\mathsf S_s,
                         \mathsf U_{\mathrm{gap}},\mathsf S_{\mathrm{gap}})
\]
over all nonempty strict linear orders
and all unary interpretations.
Conversely, every formula of this
explicitly defined temporal language
has an equivalent monadic first-order
formula over the same model class.

The forward construction is the structural
first-order induction using simple
partition atoms, the complete negation
operation, positive existential elimination,
and the full one-variable temporal return.
All canonical label names are inlined.
The converse uses the four proved
first-order truth tables. Thus both
directions are uniform and constructive
\cite[Theorems~2.3 and 4.4]{Rabinovich2018Stavi}.
This is the later defined Stavi basis;
it does not assign an unprinted literal
definition to GPSS's \(U',S'\).
\end{hbtheorem}

\begin{hbexample}{def:p09-stavi-completion-examples}{A forced point and a missing cut}
On \(\mathbb Q\), take \(a=0,b=2\), let \(g\)
be the cut at \(\sqrt2\), and put
\(D_1(x)\equiv x<g\), \(D_2=\top\),
\(D_3(x)\equiv g<x\). These are arbitrary
unary interpretations; \(g\) is not a
parameter of the first-order language.
The two failure gaps coincide. With
no flags, three cells
\((0,1),[1,3/2],(3/2,2)\) realize the
partition. Its middle cell crosses \(g\),
so \(\lambda_2(0)\land\rho_2(2)\) holds.
If the middle cell is forced to be
a singleton, no partition exists:
placing that point below \(g\) leaves
points below \(g\) in the final cell,
and placing it above \(g\) leaves
points above \(g\) in the initial cell.
The cut point itself is absent.
The equal-gap formula handles these
two outcomes by retaining or excluding
the within-cell index \(2\).

Here the test \(T_W\) uses \(D_3\)
and is false. Substituting \(D_2=\top\)
would make it true and discard the
required equal-gap case. Independently,
on \(\mathbb R\) the two-cell word
\((\top,\top)\) with the first cell
forced to be a singleton cannot
partition \((0,1)\). The formula
\(\neg X\top(0)\) detects the obstruction;
the uniform-first handler requires
its stated unflagged hypothesis.
\end{hbexample}
